\chapter{Amblyopia}
\index{Amblyopia}

\section*{Definition and Overview}
Amblyopia, commonly known as lazy eye, is reduced vision in one eye that develops during childhood because the brain does not learn to process the visual input from that eye correctly. The eye itself usually looks normal and may be perfectly healthy, but the visual brain circuits that should serve it never develop fully, so the brain increasingly ignores its image. Amblyopia is the most common cause of reduced vision in children and is highly treatable in early childhood but very difficult to reverse in later life. Early detection through screening, and treatment before the age of about seven or eight, are the keys to a good outcome. Treatment works by encouraging the brain to use the weaker eye until its vision improves.

\section*{Epidemiology}
Amblyopia affects approximately 1 to 3\% of children, making it the commonest cause of visual impairment in this age group. It is more likely to develop in children with a squint (strabismus), a large difference in spectacle prescription between the two eyes, or a condition that blocks light in one eye, such as a cataract, a droopy eyelid, or corneal opacity. Without screening it frequently goes unnoticed, because the good eye sees normally and the child rarely complains. Risk is increased in children born prematurely, those with a family history of squint or amblyopia, and those with developmental delay.

\section*{Aetiology and Pathophysiology}
Amblyopia develops when the brain does not receive clear, matching images from both eyes during the critical period of visual development in early childhood:

\begin{itemize}
    \item \textbf{Strabismic amblyopia:} when the eyes point in different directions, the brain suppresses the image from the turned eye to avoid double vision, and that eye's visual connections fail to mature.
    \item \textbf{Refractive (anisometropic) amblyopia:} a large difference in long-sightedness, short-sightedness, or astigmatism between the two eyes means the brain always prefers the clearer image.
    \item \textbf{Deprivation amblyopia:} anything that blocks light in one eye during early life, such as a congenital cataract, ptosis, or corneal scar, prevents the formation of a clear image and is the most severe form.
    \item \textbf{Critical period:} vision develops mainly in the first years of life; if one eye receives poor images during this window, the corresponding brain connections never mature properly.
    \item \textbf{Genetics:} squint and significant refractive errors run in families, increasing the risk for close relatives.
\end{itemize}

\section*{Clinical Features}
Amblyopia is often symptom-free, which is precisely why screening matters:

\begin{itemize}
    \item Poor vision in one eye that is not fully corrected by glasses alone.
    \item A squint (a turned eye), which may be constant or intermittent.
    \item Squinting, closing one eye, or tilting the head to see better.
    \item Poor depth perception and difficulty judging distances.
    \item Rubbing the affected eye or complaining of blurry vision.
    \item The child may be completely unaware that one eye sees poorly, since the good eye compensates.
\end{itemize}

\section*{Differential Diagnosis}
Reduced vision in a child can have other causes that must be excluded:

\begin{itemize}
    \item Simple refractive error that responds to glasses alone.
    \item Cataract, congenital or acquired.
    \item Retinal disease, including inherited retinal dystrophies.
    \item Optic nerve disorders.
    \item Corneal scarring or opacity.
    \item Childhood glaucoma.
    \item Retinoblastoma or another intraocular tumour, which is rare but serious.
\end{itemize}

\section*{When to Seek Medical Care}
Arrange an eye examination for any child who fails a vision screen, appears to have a turned eye, squints, or complains of poor vision. Children at higher risk, such as those born prematurely or with a family history of squint, should have early eye checks even without symptoms.

\textbf{Contact your local emergency services for:}
\begin{itemize}
    \item Sudden loss of vision in one or both eyes.
    \item A new white reflection or glow in the pupil, which can be a sign of retinoblastoma.
    \item Sudden painful redness of the eye.
    \item Sudden double vision.
    \item Any sudden change in vision that causes concern.
\end{itemize}

\section*{Diagnosis and Investigations}
The diagnosis of amblyopia is made by an orthoptist or ophthalmologist:

\begin{itemize}
    \item \textbf{Vision testing:} age-appropriate eye charts, such as picture tests for young children, to compare the two eyes.
    \item \textbf{Refraction:} eye drops are used to relax the focusing muscles so the full spectacle prescription can be measured accurately.
    \item \textbf{Squint assessment:} cover tests and assessment of eye movements and alignment.
    \item \textbf{Examination of the eye:} slit-lamp examination and inspection of the back of the eye (fundoscopy) to exclude structural causes.
    \item \textbf{Photoscreening:} automated screening devices used in community vision screening programs.
    \item \textbf{Referral:} to an ophthalmologist for confirmation of the diagnosis and ongoing monitoring of treatment.
\end{itemize}

\section*{Management}
Treatment works by forcing the brain to use the weaker eye and is most effective in early childhood:

\begin{itemize}
    \item \textbf{Glasses:} correcting any refractive error is usually the first step and may improve vision substantially on its own.
    \item \textbf{Patching (occlusion):} covering the good eye for set hours each day forces the brain to use the amblyopic eye.
    \item \textbf{Atropine drops:} blurring the vision of the good eye as an alternative to patching for children who resist occlusion.
    \item \textbf{Treatment of the underlying cause:} surgery for cataract or ptosis, and squint surgery to straighten the eyes where alignment is poor.
    \item \textbf{Ongoing monitoring:} treatment requires regular review, and patching may need to be restarted if vision slips.
    \item \textbf{Timing:} benefit diminishes with age, and treatment after about eight to ten years is usually much less effective.
\end{itemize}

\section*{Complications}
If left untreated beyond the critical window, the loss of vision is usually permanent:

\begin{itemize}
    \item Permanent reduced vision in the affected eye.
    \item Reduced depth perception and a narrowed visual field on the affected side.
    \item Reduced career options requiring good binocular vision, such as aviation, the armed forces, and emergency services.
    \item Increased risk of serious visual impairment if the good eye is later injured or diseased.
    \item Cosmetic and social concerns associated with an uncorrected squint.
\end{itemize}

\section*{Prognosis and Outcomes}
With early treatment, most children achieve useful and often good vision in the affected eye. Success depends on the age at diagnosis, the underlying cause, and how consistently treatment is followed, and improvement may take months of patching. Patching rarely restores vision completely, and the gain can regress if treatment stops too soon, so maintenance is important. After about the age of eight to ten, treatment is usually much less effective, which is why early detection through screening remains the single most important factor in a good outcome.

\section*{Prevention and Self-Care}
The most important preventive measure is early detection through screening:

\begin{itemize}
    \item Attend recommended vision screening appointments for all children.
    \item Have children with a family history of squint or eye problems checked early, even without symptoms.
    \item Keep follow-up appointments and follow the prescribed patching schedule consistently.
    \item Encourage the child to use the weaker eye during treatment with age-appropriate activities.
    \item Report any change in the child's eyes, vision, or behaviour promptly.
    \item Protect eyes from injury with appropriate safety glasses for sport.
\end{itemize}

\section*{Sources and Further Reading}
The following reputable sources provide reliable, up-to-date information on amblyopia:

\begin{itemize}
    \item NHS. \textit{Lazy eye (amblyopia) information for patients.} https://www.nhs.uk/
    \item Mayo Clinic. \textit{Lazy eye (amblyopia): symptoms, causes, and treatment.} https://www.mayoclinic.org/
    \item MSD Manual. \textit{Amblyopia overview.} https://www.msdmanuals.com/
    \item MedlinePlus. \textit{Amblyopia health information.} https://medlineplus.gov/
    \item World Health Organization. \textit{Information on blindness and vision impairment prevention.} https://www.who.int/
    \item NICE. \textit{Clinical guidance on eye disorders in children.} https://www.nice.org.uk/
\end{itemize}
